We asked when a learned correction should defer to a friction model at breakaway, and turned the answer into shields with closed-form certificates for pivoting and sliding. The answer depends on the task's physics more than on the learner. In sliding, a certified action almost always exists. A residual trained under the shield then keeps overshoot near zero across friction ratios, and degrades gradually when $r_{\max}$ is underestimated, although it gives up most of the unshielded residual's measurement saving. In pivoting, the stick--slip energy budget leaves no certified release for most shallow targets. There a shield cannot make a learned correction both efficient and safe: it can only choose between the model's overshoot and stopping short, and the feasibility condition $r<\varphi_{\lim}/\sin\varphi_{\lim}$ predicts in advance which targets force that choice. We see two next steps: a fast slip-triggered catch, which would enlarge the feasible region for pivoting, and measuring $r_{\max}$ on real pad--object pairs to test whether these certificates survive contact with hardware.
